\documentclass[11pt, a4paper]{report}
\usepackage[spanish]{babel}
\usepackage{geometry}
\usepackage{graphicx}
\usepackage{amsmath}
\usepackage{apacite}
\usepackage{tabularx}
\usepackage{listings}
\usepackage{xcolor}
\usepackage{eso-pic}
\usepackage{setspace}
\usepackage{titlesec}
\usepackage{fancyhdr}

% --- CONFIGURACIÓN DE FUENTE (Igual al TFG) ---
\usepackage[scaled]{helvet}
\renewcommand*\familydefault{\sfdefault}
\usepackage[T1]{fontenc}

\geometry{a4paper, margin=2.5cm, includeheadfoot}
\setstretch{1.2}
\setlength{\parskip}{0.6cm}
\setlength{\parindent}{0pt}

% Configuración de títulos (Máximo 16pt)
\titleformat{\chapter}[display]
  {\normalfont\Large\bfseries}{\chaptertitlename\ \thechapter}{10pt}{\Large}
\titleformat{\section}
  {\normalfont\large\bfseries}{\thesection}{1em}{}
\titleformat{\subsection}
  {\normalfont\normalsize\bfseries}{\thesubsection}{1em}{}
\titlespacing*{\chapter}{0pt}{0pt}{20pt}

% Configuración de encabezado y pie de página
\pagestyle{fancy}
\fancyhf{}
\fancyhead[L]{Sistemas Inteligentes}
\fancyhead[R]{Teoría}
\fancyfoot[C]{\thepage}
\renewcommand{\headrulewidth}{0.4pt}
\renewcommand{\footrulewidth}{0.4pt}

\title{%
    \vspace*{2cm}
    {\Large UNIVERSIDAD DE HUELVA}\\[0.5cm] % 
    {\huge \textbf{Sistemas Inteligentes}}\\[0.5cm]
    {\Large Teoría}\\[2cm]
}


\author{%
    \vspace{1cm}
    Escuela Tecnica Superior De Ingeniería\\[0.5cm]
    Departamento De Tecnologias De La Información \\[0.5cm]
}

\date{%
    \vspace{2cm}
    \large \today
}

\lstset{
    language=Java,
    basicstyle=\small\ttfamily,
    keywordstyle=\color{blue}\bfseries,
    commentstyle=\color{green!50!black}\itshape,
    stringstyle=\color{orange},
    numbers=left,
    numberstyle=\tiny\color{gray},
    breaklines=true,
    frame=single,
    tabsize=2,
    captionpos=b
}

% Cambiamos el nombre de los capítulos a "Tema"
\addto\captionsspanish{\renewcommand{\chaptername}{Tema}}
\begin{document}
\input{Portada.tex}

% A partir de aquí, el documento vuelve a la normalidad (sin fondo)
\tableofcontents
\newpage


% Capitulos del libro
\chapter{Introducción: Búsquedas en IA}

\section{Formulación}
La Inteligencia Artificial Clásica trata entornos con las siguientes características:
\begin{itemize}
    \item Estático: El entorno no varía.
    \item Observable: Conocemos todas las variables que nos interesan.
    \item Discreto: Los valores tienen límite de división.
    \item Determinístico: El estado siguiente es consecuencia del estado actual y de la acción realizada.
\end{itemize}

\section{Características}
Se puede construir un grafo con la representación del problema. Las características de este grafo serán:
\begin{itemize}
    \item Cada estado se representará mediante un nodo.
    \item Deberá tener un estado inicial y, al menos, un estado final.
    \item Las aristas representan las acciones entre los estados.
    \item Puede ser cíclico.
\end{itemize}

\section{Espacio De Estados}
\subsection{Árbol De Búsqueda}

\section{Solución}

\section{Ejemplos}

\section{Búsqueda De Soluciones}

\section{Algoritmo General De Búsqueda}

\section{Estrategias De Búsqueda}

\section{Criterios Para Evaluar Las Estrategias}

\section{Búsquedas Con Adversarios}



\section{Búsqueda En El Espacio De Estados (Temp)}


\subsection{Estructuras de Datos y Definiciones}
A continuación se define la estructura necesaria para un algoritmo de búsqueda genérico (como A* o Dijkstra).

\begin{lstlisting}[caption=Interfaces y Clase Nodo]
public interface Estado {
    public Estado aplicarAccion(Action a);
    public boolean esMeta();
    public boolean esValido();
    public int h(); // Heuristica: coste estimado hasta la meta
}

public interface Action {
    public int costeAccion();
}

public class Nodo implements Comparable<Nodo> {
    private int g;            // Coste acumulado (inicio -> actual)
    private Estado estado;
    private Nodo padre;
    private Action accion;

    public Nodo(int g, Estado estado, Nodo padre, Action accion) {
        this.g = g;
        this.estado = estado;
        this.padre = padre;
        this.accion = accion;
    }

    public int getF() { return this.g + this.estado.h(); }
    public int getG() { return g; }
    public Estado getEstado() { return estado; }
    public Nodo getPadre() { return padre; }
    public Action getAccion() { return accion; }

    @Override
    public int compareTo(Nodo otro) {
        return Integer.compare(this.getF(), otro.getF());
    }
}
\end{lstlisting}

\subsection{Algoritmo de Búsqueda General}
El siguiente pseudocódigo implementa una búsqueda que evita ciclos mediante un conjunto de estados cerrados y utiliza una cola de prioridad para optimizar la expansión.

\begin{lstlisting}[caption=Algoritmo de Búsqueda corregido]
public class BusquedaEspacioEstados {
    public List<Action> buscar(Estado inicial, List<Action> acciones) {
        PriorityQueue<Nodo> abiertos = new PriorityQueue<>();
        Set<Estado> cerrados = new HashSet<>();

        abiertos.add(new Nodo(0, inicial, null, null));

        while (!abiertos.isEmpty()) {
            Nodo actual = abiertos.poll();
            Estado estadoActual = actual.getEstado();

            if (estadoActual.esMeta()) {
                return reconstruirCamino(actual);
            }

            if (!cerrados.contains(estadoActual)) {
                cerrados.add(estadoActual);

                for (Action a : acciones) {
                    Estado sigEstado = estadoActual.aplicarAccion(a);
                    if (sigEstado != null && sigEstado.esValido()) {
                        int nuevoG = actual.getG() + a.costeAccion();
                        abiertos.add(new Nodo(nuevoG, sigEstado, actual, a));
                    }
                }
            }
        }
        return null; // No hay solucion
    }

    private List<Action> reconstruirCamino(Nodo nodo) {
        List<Action> camino = new ArrayList<>();
        Nodo aux = nodo;
        while (aux != null && aux.getAccion() != null) {
            camino.add(0, aux.getAccion());
            aux = aux.getPadre();
        }
        return camino;
    }
}
\end{lstlisting}

\subsection{Notas de Implementación}
\begin{itemize}
    \item \textbf{Variantes del algoritmo}:
    \begin{itemize}
        \item \textbf{A*}: Se activa cuando \texttt{h()} devuelve una estimación optimista del coste restante.
        \item \textbf{Dijkstra}: Se obtiene si la función heurística \texttt{h()} devuelve siempre 0.
        \item \textbf{BFS}: Se obtiene si el coste de toda acción es 1 y \texttt{h()} es 0.
    \end{itemize}
    \item \textbf{Control de Ciclos}: El uso de \texttt{Set<Estado> cerrados} es vital para evitar bucles infinitos en grafos.
    \item \textbf{Requisitos de Java}: Para que la detección de estados repetidos funcione correctamente, la clase que implemente \texttt{Estado} debe sobreescribir obligatoriamente los métodos \texttt{equals()} y \texttt{hashCode()}.
\end{itemize}

% ---------Bibliografía del Tema------------
\bigskip
\noindent\rule{\textwidth}{0.4pt} 

\section*{Bibliografía del Tema}
\addcontentsline{toc}{section}{Bibliografía del Tema}

\begin{itemize}
    \item \textbf{Russell, S. \& Norvig, P.} (2003). \textit{Artificial Intelligence: A Modern Approach} (2nd ed.). Prentice Hall. Chapter 3.
    
\end{itemize}
\newpage
\chapter{Agentes Inteligentes}
\section{Conceptos Fundamentales}

Para comprender la materia, el primer paso fundamental es definir qué es un agente. Partiendo de una base lingüística, la \textbf{Real Academia Española (RAE)} propone varias acepciones que resultan relevantes:

\begin{quote}
\textit{Agente:} 
\begin{enumerate}
    \item adj. Que obra o tiene virtud de obrar.
    \item adj. Gram. Dicho de una palabra o de una expresión: Que designa a la persona, animal o cosa que realiza la acción del verbo.
    \item m. Persona o cosa que produce un efecto. 
    \item m. Persona que obra con poder de otra.
\end{enumerate}
\end{quote}


Sin embargo, al trasladar este concepto al ámbito de la informática, nos encontramos con un problema: \textbf{no existe una definición de agente software universalmente aceptada}. La terminología es diversa y la definición suele depender del dominio de aplicación o del autor consultado.

\medskip
Nosotros vamos a considerar que un agente es cualquier entidad capaz de percibir su entorno mediante \textbf{sensores} y actuar sobre él a través de \textbf{efectores}. Aunque existen múltiples definiciones en la literatura, se acepta que un agente es un sisema situado que busca cumplir sus objetivos a lo largo del tiempo.

\medskip
Un agente basicamente hace:

\begin{itemize}
    \item Percibe el entorno.
    \item Actúa sobre el entorno.
    \item Asigna perciones a acciones.
    \item Mide el rendimiento.
\end{itemize}

Es decir, es todo aquello que percibe su ambiente mediante sensores y responde o actúa en el ambiente por medio de efectores. No requieren participación de los humanos para operar.

\section{Percepciones y Acciones}
El comportamiento de un agente depende de su secuencia de percepciones en un momento dado.

\subsection{El ciclo de percepción-acción}
Un agente ideal se define por un mapeo que especifica qué acción debe emprender ante cualquier secuencia de percepciones. Matemáticamente:
\begin{equation*}
    f: \mathcal{P}^* \to \mathcal{A}
\end{equation*}

Donde:
\begin{itemize}
    \item $\mathcal{P}^*$ representa el historial o secuencia de todas las percepciones recibidas hasta el momento.
    \item $\mathcal{A}$ es la acción resultante elegida por el agente.
\end{itemize}

\subsection{Sensores y Efectores}
\begin{itemize}
    \item \textbf{Sensores:} Adquieren información del medio ambiente. Pueden presentar errores o incertidumbre debido a cambios en el medio. Ejemplos: cámaras, acelerómetros, sonares.
    \item \textbf{Efectores:} Permiten al agente actuar. En humanos son los músculos y en máquinas son motores, válvulas o brazos mecánicos.
\end{itemize}

\section{Comportamiento}


El comportamiento de un agente depende de la secuencia de percepciones en un momento dado. Se puede caracterizar un agente elaborando una lista de percepciones-acciones. A esto se le conoce como mapeo de secuencias de percepciones para acciones.

\medskip
El mapeo ideal es especificar qué tipo de acción deberá emprender un agente como respuesta a una determinada secuencia de percepciones.


% ---------------------------------------------------------------------------------------
% ---------------------------------------------------------------------------------------
% ---------------------------------------------------------------------------------------
\section{Agentificacion}
Antes de diseñar un programa de agente tenemos que hacer la descripción de las percepciones, las acciones, las metas y del ambiente.
La arquitectura de un agente es la siguiente:
\begin{equation*}
    \text{Agente} = \text{Arquitectura} + \text{Programa}
\end{equation*}

Donde:
\begin{itemize}
    \item \textbf{Arquitectura:} Pone al alcance del programa las percepciones obtenidas mediante los sensores, lo ejecuta y alimenta el actuador con acciones elegidas por el programa conforme se van generando.
    \item \textbf{Programa:} Es un algoritmo que recibe las percepciones del agente y genera una secuencia de acciones.
\end{itemize}

\section{Agentes Inteligentes}
Los agentes inteligentes tienen las siguientes características:
\begin{itemize}
    \item \textbf{Autonomía}: Por la propia naturaleza de los agentes, se espera que sean capaces de operar sin intervención humana directa. Esto implica que deben ser capaces de tomar decisiones y actuar de manera independiente para alcanzar sus objetivos.
    \item \textbf{Reactividad}: Los agentes inteligentes deben ser capaces de percibir su entorno y responder de manera oportuna a los cambios que ocurren en él. Esto les permite adaptarse a situaciones imprevistas y mantener un comportamiento coherente con sus objetivos.
    \item \textbf{Proactividad}: Los agentes son capaces de tomar la iniciativa para la consecución de objetivos. Esto significa que no solo reaccionan a los estímulos del entorno, sino que también pueden planificar y ejecutar acciones para alcanzar metas específicas.
    \item \textbf{Habilidad Social}: Los agentes interactúan con otros agentes (hasta humanos) a través de algún lenguaje. Esto les permite colaborar, negociar y compartir información para lograr objetivos comunes o resolver conflictos.
\end{itemize}

\subsection{Agentes Inteligentes Racionales}
Para cada posible secuencia de percepciones, un agente racional seleccionaría un acción que
maximice su medida de rendimiento, a partir de esa secuencia de acciones y el conocimiento 
que tenga en su interior. Esto es lo que se conoce como Agente Racional, y dependen de:

\begin{itemize}
    \item Medida Del Grado De Éxito.
    \item Secuencia De Percepciones.
    \item Conocimiento Acerca Del Medio.
    \item Acciones Que Puede Emprender.
\end{itemize}

\subsection{Medida de Rendimiento y Racionalidad}

Para determinar la eficacia de un agente, es imprescindible establecer un criterio de evaluación denominado \textbf{medida de rendimiento} o medida del desempeño. Este concepto se fundamenta en los siguientes pilares:

\begin{itemize}
    \item \textbf{Evaluación cualitativa:} Analiza el ``cómo'' el agente interactúa con su entorno.
    \item \textbf{Grado de éxito:} Determina qué tan exitoso ha sido un agente en el cumplimiento de su tarea.
    \item \textbf{Objetividad:} La medida debe ser objetiva, basándose en criterios externos y no en el estado interno del agente.
\end{itemize}

\bigskip

Es fundamental distinguir la racionalidad de otros conceptos cognitivos superiores. 
Según la literatura de la asignatura, la racionalidad \textbf{NO ES} omnisciencia 
(saberlo todo), clarividencia (predecir el futuro) ni implica ser exitosa necesariamente 
en todos los casos físicos.

\bigskip

En cambio, la racionalidad se puede definir como un \textbf{éxito esperado}, tomando como 
base estrictamente lo que el agente ha percibido hasta el momento a través de sus sensores. 
Por tanto, un agente es racional si elige la mejor acción posible dada la información 
limitada de la que dispone y su conocimiento previo.


\section{Entorno}
Antes de proceder al diseño de un programa de agente, es imperativo realizar una descripción exhaustiva del ambiente en el que este operará. El éxito de un agente racional depende de cómo sus sensores captan la información del entorno y cómo sus acciones lo modifican. 

\medskip
\noindent
Para categorizar estos escenarios, utilizamos una taxonomía basada en cinco dimensiones fundamentales que determinan la complejidad del problema a resolver:

\bigskip

\begin{table}[h]
\centering
\renewcommand{\arraystretch}{1.5}
\begin{tabularx}{\textwidth}{|l|X|} 
\hline
\textbf{Propiedad} & \textbf{Descripción} \\ \hline
\textbf{Accesibilidad} & Es \textbf{accesible} si los sensores detectan todos los aspectos que el agente requiere para elegir una acción. De lo contrario, es no accesible. \\ \hline
\textbf{Determinismo} & Es \textbf{determinista} si el estado siguiente del entorno se puede determinar completamente por el estado actual y las acciones del agente. \\ \hline
\textbf{Episodicidad} & Es \textbf{episódico} cuando la experiencia del agente se divide en episodios independientes; cada acción corresponde solo a las percepciones de su episodio. \\ \hline
\textbf{Dinamismo} & Es \textbf{estático} si el ambiente no cambia mientras el agente se encuentra decidiendo la acción. Si cambia, se considera dinámico \\ \hline
\textbf{Continuidad} & Es \textbf{discreto} si existe una cantidad limitada y distinguible de percepciones y acciones distintas. \\ \hline
\end{tabularx}
\caption{Dimensiones y características de los entornos para agentes inteligentes.}
\end{table}

\section{Tipos De Agentes}
Antes de profundizar en los sistemas computacionales, es necesario situar al agente dentro de una taxonomía general basada en su naturaleza y el entorno en el que opera. En este sentido, podemos distinguir tres categorías fundamentales:
\begin{itemize}
    \item \textbf{Agentes Naturales:} Aquellos constituidos por cuerpos biológicos que operan en entornos naturales, donde su medida de rendimiento suele estar ligada a la supervivencia o la reproducción.
    
    \item \textbf{Agentes de Hardware (Robots):} Sistemas artificiales que actúan directamente en el entorno físico mediante el uso de sensores como cámaras u odómetros, y efectores como ruedas o brazos mecánicos.
    
    \item \textbf{Agentes Software (Softbots):} Son programas que habitan en entornos virtuales, como sistemas operativos o Internet. A diferencia de los robots físicos, sus sensores y efectores son componentes lógicos dependientes del dominio, como protocolos de red o consultas a bases de datos.
\end{itemize}

\bigskip

De aqui en adelante, nos centraremos específicamente en el estudio de los \textbf{Agentes Software}. Estos sistemas son fundamentales en áreas como el comercio electrónico, la recuperación de información y la gestión del conocimiento, donde su capacidad para automatizar tareas complejas de forma autónoma resulta indispensable.

\subsection{Reactivos} Selecciona una acción únicamente en función de la percepción actual. Se implementa con reglas condición-acción.

\subsection{Reactivos Basados En Modelos} Se utilizan para abordar entornos parcialmente observables. Mantiene un estado interno que se actualiza
a lo largo del tiempo usando el conocimiento adquirido. Usa un modelo del mundo y un conjunto de reglas.

\subsection{Basados En Objetivos} Buscan lograr ciertos objetivos. El comportamiento se complica cuando las secuencias de acciones son
largas. Para ello se utilizan estrategias de búsqueda y planificación. Se tiene en cuenta el futuro.

\subsection{Basados En Utilidad} Ciertas metas pueden ser alcanzadas de diferentes formas. La función de utilidad mapea una secuencia
de estados en un número real, la utilidad. La utilidad es una función que correlaciona un estado y un número
real mediante el que se caracteriza el grado de satisfacción.

\subsection{Agentes Aprendices} El agente tiene un elemento de aprendizaje que mejora el rendimiento de este. Nos provee feedback sobre
el rendimiento del agente. La elección de la acción se realiza en función de las percepciones y del elemento
de aprendizaje.


\bigskip
\noindent\rule{\textwidth}{0.4pt} % Línea divisoria para separar del contenido

\section*{Bibliografía del Tema}
\addcontentsline{toc}{section}{Bibliografía del Tema}

\begin{itemize}
    \item \textbf{Russell, S. \& Norvig, P.} (2003). \textit{Artificial Intelligence: A Modern Approach} (2nd ed.). Prentice Hall. Chapter 2.
    
    \item \textbf{Wooldridge, M.} (2002). \textit{An Introduction to Multiagent Systems}. Wiley. Chapters 1 and 2.
\end{itemize}

\newpage


\chapter{Comportamientos}

\section{Introducción A Los Comportamientos Inteligentes}
\subsection{Niveles De IA En Agentes}
Se distinguen tres niveles principales en el desarrollo de agentes inteligentes:
\begin{itemize}
    \item \textbf{Básico}: Basado en técnicas de toma de decisiones principalmente reactivas.
    \item \textbf{Avanzado}: Incluye técnicas de razonamiento y planificación.
    \item \textbf{Futuro}: Temas de investigación actualmente en desarrollo.
\end{itemize}

\subsection{Tipos De Algoritmos}
\subsubsection{Sin Búsqueda}
En estos algoritmos, el coste computacional es predecible. El conocimiento reside en la lógica del código o en tablas externas. Ejemplos comunes son los sistemas de reglas y los árboles de decisión.

\subsubsection{Con Búsqueda}
El coste computacional depende del tamaño del espacio de búsqueda, el cual suele ser grande. El conocimiento se encuentra en la evaluación de los estados. Incluye algoritmos como Minimax, Planificación y Pathfinding.


\section{Sistemas Basados En Reglas}
También conocidos como sistemas de producción o sistemas expertos. Los sistemas basados en reglas son uno de los paradigmas de IA más
exitosos. Se originaron en los años 70 y 80.

\subsection{Arquitectura Y Componentes}
Los datos son alamcenados como hechos y reglas. Sus componenetes son:
\begin{itemize}
    \item \textbf{Memoria de trabajo}: Almacena los hechos ciertos o el modelo del mundo.
    \item \textbf{Reglas de producción}: Definidas habitualmente como estructuras IF...THEN.
    \item \textbf{Motor de inferencia}: Programa encargado de calcular las consecuencias y aplicar las reglas.
    \item \textbf{Resolución de conflictos}: Decide qué regla aplicar si más de una es válida.
\end{itemize}

\subsection{Mecanismos De Razonamiento}
Una regla se ``dispara'' cuando sus antecedentes coinciden con los hechos de la base de datos.  El
disparo de la regla añade a la base de hechos el consecuente de esta. 
\medskip
Existen dos enfoques de inferencia en los sistemas basados en reglas:

\subsubsection{Encadenamiento Hacia Delante}
Tambien llamado ``Forward Chaining'', es un razonamiento deductivo basado en datos. Parte de antecedentes para llegar a una conclusión o acción.

\subsubsection{Encadenamiento Hacia Atrás}
Tambien llamado ``Backward Chaining'', es un razonamiento inductivo impulsado por objetivos. Parte de una conclusión e intenta demostrarla buscando los antecedentes necesarios en la base de hechos.

\subsection{Estrategis De Resolución De Conflictos}
Existen varias formas de solucionar los conflictos:
\begin{itemize}
    \item \textbf{Orden físico}: Las reglas se evalúan según su orden en el código.  Es difícil añadir reglas a estos sistemas.
    \item \textbf{Prioridad de datos}: Se utiliza una cola de prioridad para los elementos del problema. Se usa la regla que ocupa los elementos más prioritarios.
    \item \textbf{Especificidad O Máxima Especialidad}: Se elige la regla con mayor número de antecedentes coincidentes.
    \item \textbf{Selección aleatoria}: Disparo aleatorio o de todas las reglas aplicables.
    \item \textbf{Lanzar Todo}: Se Disparan todas las reglas aplicables.
    \item \textbf{Mejor regla}: Se busca la regla más adecuada para lanzarla.
\end{itemize}

\subsection{Ventajas Y Desvantajas}
\begin{itemize}
    \item \textbf{Ventajas}: Son universales, muy expresivos, modulares y fáciles de leer.
    \item \textbf{Desventajas}: No todo el conocimiento encaja en reglas, requieren consenso de expertos, pueden ser intensivos en memoria/computación, pueden ser dificiles de depurar, requiere un conocimiento detallado y es dificil olvidad hechos.
\end{itemize}

\section{Árboles De Decisión}
\subsection{Carasterísticas}
Se consideran la técnica de IA más básica. Destacan por ser fáciles de implementar, rápidos de ejecutar y sencillos de entender.

\subsection{Implementación}
Se utiliza Programación Orientada a Objetos (POO) para evitar estructuras de código frágiles. Los nodos interiores representan decisiones y las hojas representan acciones.

\subsection{Funcionamiento}
El test de cada nodo suele ser de tiempo constante. El tiempo total de ejecución depende de la profundidad del árbol. Se recomienda construir árboles balanceados y colocar las decisiones más costosas al final.

\subsection{Ventajas Y Desvantajas}
\begin{itemize}
    \item \textbf{Ventajas}: Son extremadamente rápidos de ejecutar ($O(depth)$), fáciles de implementar y depurar, y muy visuales para diseñadores no programadores.
    \item \textbf{Desventajas}: Pueden volverse inmanejables si el árbol es muy profundo, presentan dificultades para gestionar estados que dependen del tiempo (no tienen memoria) y suelen conllevar duplicidad de comprobaciones en distintas ramas.
\end{itemize}

\section{Máquinas De Estado Finitas (FMS)}
\subsection{Carasterísticas}
Permiten una correspondencia natural entre estados y comportamientos. Formalmente constan de un conjunto de estados, un estado inicial, un alfabeto de entrada y transiciones.

\subsection{Implementación}
\begin{itemize}
    \item \textbf{Inicial}: Basada en enumerados y sentencias condicionales; eficiente pero difícil de mantener.
    \item \textbf{Orientada a Objetos}: Clases para \texttt{State} y \texttt{Transition}; más limpia y flexible.
\end{itemize}

\subsection{Combinar Árboles De Decisión Con Máquinas De Estados}
Se utilizan para evitar comprobaciones duplicadas o costosas en las transiciones de la FSM. Por ejemplo, si una condición como "jugador a la vista" es computacionalmente cara, el árbol de decisión optimiza la evaluación.

\subsection{Ventajas Y Desvantajas}
\begin{itemize}
    \item \textbf{Ventajas}: Son computacionalmente económicas, fáciles de diseñar para comportamientos simples y permiten una transición clara y predecible entre estados lógicos.
    \item \textbf{Desventajas}: Sufren de la "explosión de estados" (el número de transiciones crece exponencialmente al añadir nuevos estados), son difíciles de reutilizar y no gestionan bien las interrupciones o tareas paralelas.
\end{itemize}

\section{Máquinas De Estado Jerárquicas  (HFSM)}
Surgieron para gestionar interrupciones (como recargar batería) sin perder el progreso del comportamiento actual.

\subsection{Ventajas Y Desvantajas}
\begin{itemize}
    \item \textbf{Ventajas}:
        \begin{itemize}
            \item \textbf{Memoria de estado}: Permiten recordar el estado interno y continuar tras una interrupción.
            \item \textbf{Reducción de la complejidad}: Evitan la explosión de estados. En un ejemplo con interrupciones de batería y enemigos, una HFSM requiere 7 estados frente a los 12 de una FSM plana.
        \end{itemize}
    \item \textbf{Desventajas}: Su implementación es significativamente más compleja que una FSM plana, requieren un mayor diseño previo para definir las jerarquías y tienen una ligera sobrecarga (overhead) adicional en el motor de inferencia.
\end{itemize}

% ---------Bibliografía del Tema------------
\bigskip
\noindent\rule{\textwidth}{0.4pt} 

\section*{Bibliografía del Tema}
\addcontentsline{toc}{section}{Bibliografía del Tema}

\begin{itemize}
    \item Buckland, M. (2005). \textit{Programming Game AI by Example}. Wordware Publishing, Inc.
    
    \item Millington, I., \& Funge, J. (2018). \textit{Artificial Intelligence for Games} (3rd ed.). CRC Press.
    
    \item Rabin, S. (Ed.). (2017). \textit{Game AI Pro 3: Collected Wisdom of Game AI Professionals}. CRC Press.
    
    \item Russell, S. J., \& Norvig, P. (2020). \textit{Artificial Intelligence: A Modern Approach} (4th ed.). Pearson.
    
    \item Schwab, B. (2008). \textit{AI Game Development: Digital Media Academy Series}. Charles River Media.
\end{itemize}
\newpage
\chapter{Planificación Clásica}

\section{Definición De Planifiación}
La planificación es el proceso de búsqueda y articulación de una secuencia de acciones (denominada \textbf{plan}) que permite alcanzar uno o varios objetivos partiendo de un estado inicial. 

Los elementos fundamentales que intervienen en un problema de planificación son:
\begin{itemize}
    \item \textbf{Estados:} Representaciones del mundo en un momento dado mediante predicados lógicos.
    \item \textbf{Acciones/Operadores:} Generan descripciones de estados nuevos al ser aplicados.
    \item \textbf{Metas (Objetivos):} El estado final o condiciones que se desean conseguir.
    \item \textbf{Plan:} Secuencia ordenada de acciones que transforman el estado inicial en el estado meta.
\end{itemize}

\section{Planificación Lineal: STRIPS}
\textbf{STRIPS} (\textit{Stanford Research Institute Problem Solver}) es tanto un lenguaje como un algoritmo de planificación que utiliza una pila de objetivos. Realiza una exploración en profundidad del espacio de estados.


\subsection{Lenguaje y Operadores}

El vocabulario que vamos a usar para definir un problema es el siguiente:
\begin{itemize}
    \item P: significa insertar P.
    \item ¬P: Significa eliminar P.
    \item Los estados se representan con un conjunto de predicados.
\end{itemize}

\bigskip
STRIPS utiliza la \textbf{Hipótesis del Mundo Cerrado}: cualquier condición no mencionada explícitamente en el estado se considera falsa. Los operadores constan de:
\begin{enumerate}
    \item \textbf{Precondiciones (P):} Literales que deben ser ciertos para ejecutar la acción.
    \item \textbf{Lista de Adición (A):} Literales que pasan a ser verdaderos tras la acción.
    \item \textbf{Lista de Supresión (S):} Literales que dejan de ser ciertos tras la acción.
\end{enumerate}

\subsection{Algoritmo de Implementación}
El proceso se basa en el manejo de una pila:
\begin{itemize}
    \item Si la cima es un \textbf{operador}: se ejecuta si sus precondiciones se cumplen; si no, se apilan sus precondiciones.
    \item Si la cima es una \textbf{meta}: si ya es cierta en el estado actual, se elimina; si no, se busca un operador que la añada y se apila dicho operador.
    \item Si la cima es una \textbf{conjunción de metas}: se elimina si todas son ciertas; de lo contrario, se apilan las metas individuales en distintas combinaciones.
\end{itemize}

\subsection{Anomalía de Sussman}
Es un problema intrínseco a la planificación lineal que ocurre cuando la consecución de un objetivo implica necesariamente destruir un objetivo alcanzado previamente. Esto sucede porque el planificador lineal asume erróneamente que los objetivos son independientes y pueden resolverse secuencialmente sin interferencias.

---
<PDDL es el lenguaje con el que vamos a escribir la planificacion
Esta escrito el logica de primer orden, los estados y los objetivos se escriben como un conjunto de formulas logicas, acciones tienen ( nombre, precondiciones, conjunto de supresin )>
<Se hizo ejemplo de strips> 

Anomalía de Sussman 
Cuando los objetivos no son independientes hay veces que el problema no tiene solucion
(No siempre aparece, podemos incluso poner los operadores en un orden concreto y no se produce)


\section{Planificación de Orden Parcial (POP)}
A diferencia de la planificación lineal, la \textbf{Planificación de Orden Parcial} (o no lineal) permite que los objetivos no tengan que seguir una secuencia única estricta desde el inicio. Se basa en el **Espacio de Planes**, donde cada nodo es un plan parcial.

\subsection{Características Principales}
\begin{itemize}
    \item \textbf{Principio de Mínimo Compromiso (Least Commitment):} No se determina el orden de las acciones a menos que sea estrictamente necesario para evitar conflictos (\textit{amenazas}). "No ordenar lo que no necesita ser ordenado".
    \item \textbf{Flexibilidad:} Permite la ejecución en paralelo de acciones independientes.
    \item \textbf{Búsqueda:} A diferencia de STRIPS, POP busca en el espacio de planes parciales, no en el espacio de estados.
\end{itemize}

\subsection{Componentes Formales de un Plan}
Un plan parcial se define como una tupla $\pi = (A, <, \mathcal{CL}, \text{Agenda})$, donde:
\begin{itemize}
    \item \textbf{Acciones (A):} Conjunto de pasos del plan. Incluye dos acciones ficticias:
    \begin{itemize}
        \item \textbf{Start:} No tiene precondiciones; sus efectos son el estado inicial.
        \item \textbf{Finish:} Sus precondiciones son los objetivos finales; no tiene efectos.
    \end{itemize}
    \item \textbf{Restricciones de Orden ($<$):} Un conjunto de pares $(A_i < A_j)$ que indican que $A_i$ debe ejecutarse antes que $A_j$.
    \item \textbf{Enlaces Causales ($\mathcal{CL}$):} Se representan como $A_i \xrightarrow{p} A_j$, significando que el efecto $p$ de la acción $A_i$ satisface la precondición $p$ de la acción $A_j$.
    \item \textbf{Agenda:} Conjunto de precondiciones abiertas (objetivos que aún no tienen un enlace causal que los satisfaga).
\end{itemize}

\subsection{Amenazas y Resolución de Conflictos}
Existe una \textbf{amenaza} cuando una acción $A_k$ tiene un efecto que niega la precondición $p$ de un enlace causal $A_i \xrightarrow{p} A_j$. Formalmente, $A_k$ es una amenaza si:
\begin{enumerate}
    \item $A_k$ tiene como efecto $\neg p$.
    \item El orden $(A_i < A_k < A_j)$ es consistente (es decir, no crea ciclos).
\end{enumerate}

Para resolverla, se añaden restricciones de orden:
\begin{itemize}
    \item \textbf{Promoción:} Se fuerza $A_j < A_k$ (la acción conflictiva va después).
    \item \textbf{Degradación:} Se fuerza $A_k < A_i$ (la acción conflictiva va antes).
\end{itemize}

\subsection{Algoritmo POP (Pseudocódigo)}
El algoritmo funciona de manera recursiva tratando de cerrar precondiciones abiertas y resolver amenazas:

\begin{enumerate}
    \item \textbf{Inicialización:} Plan con $\{\text{Start, Finish}\}$, orden $\text{Start} < \text{Finish}$, y agenda con las precondiciones de \text{Finish}.
    \item \textbf{Selección:} Elegir una precondición abierta $(p, A_{next})$ de la agenda.
    \item \textbf{Selección de Operador:} Elegir una acción $A_{prev}$ (nueva o existente) que tenga a $p$ como efecto.
    \item \textbf{Actualización:}
    \begin{itemize}
        \item Añadir el enlace causal $A_{prev} \xrightarrow{p} A_{next}$.
        \item Añadir la restricción de orden $A_{prev} < A_{next}$.
        \item Si $A_{prev}$ es nueva, añadirla al plan y poner sus precondiciones en la agenda.
    \end{itemize}
    \item \textbf{Resolución de Amenazas:} Por cada acción $A_k$ que amenace un enlace, intentar promoción o degradación. Si ambas causan un ciclo, hacer \textit{backtracking}.
    \item \textbf{Bucle:} Repetir hasta que la agenda esté vacía.
\end{enumerate}

---
..-Se dice plan parcial abierto si existe alguna precondicion que no este satisfecha
Por otro lado, se dice que el plan parcial esta cerrado si todas las precondiciones estan satisfechas o vacias.
El objetivo es pasar de un plan parcial abierto a uno cerrado.

Hay dos formas de añadir reglas, por promocion y por degradación. (Solucion de amenazas)

Hace falta una serie de comprimisos para hacer pop:
\begin{itemize}
    \item Principio de minimo Compromiso
    \item El planificador es libre de añadir todas las acciones cuando y como quiera. Puede haber acciones repetidas (Mismo nombre).
    \item 
\end{itemize}
<>Hay que usar PDDL, añadir al resumen



\section{Otras Técnicas de Planificación}
\begin{itemize}
    \item \textbf{Descomposición Jerárquica (HTN):} Divide tareas complejas en subtareas más simples. Utiliza macrooperadores o niveles de abstracción para manejar planes extensos.
    \item \textbf{Planificación Clásica:} Se aplica en entornos observables, deterministas, finitos, estáticos y discretos.
\end{itemize}



% ---------Bibliografía del Tema------------
\bigskip
\noindent\rule{\textwidth}{0.4pt} 

\section*{Bibliografía del Tema}
\addcontentsline{toc}{section}{Bibliografía del Tema}

\begin{itemize}
    \item
\end{itemize}
\newpage
\chapter{Programación de Restricciones}

% Satisfaccion de restricciones
Los problemas se componen por tres elementos:
\begin{itemize}
    \item Variables
    \item Dominio de las Variables
    \item Restricciones
\end{itemize}

Buscamos en profundidad porque buscamos una solucion, todas las soluciones van a tener la misma longitud. 
La anchura va creciendo en cada nivel por lo que tardaria mas a pesar de que daria todas las soluciones incluyendo la optima si existe, en principio no hay una optima.

% Euristicas
Primero lo que mas va a fallar:
- Una forma de optimizar los algoritmos es ir buscando las opciones que antes van a fallar para agilizar el arbol de busqueda un monton, por ejemplo buscar las que menos dominio tienen primero.
- Tambien podemos mirar la que tiene mayor dominio.
- La que participe en mas restricciones.
- A la hora de chequear las restricciones no lo mismo chequear la que este en restricciones mas pesadas que en lass que estan en mas livianas.
- ordenar las restricciones  de tal forma que se ordenen de forma que se comprueben primero las restricciones gordas, esto es al generar sucesores.

Tecnica fordward chaining: (o fordward cheking, ha dicho ambos nombres no se cual es) quito todos los valores de las otras variables incompatibles con el padre, pero solo a los hijos de un cierto paso.



\section{Algoritmos}

Algoritmo de Satisfaccion de restricciones: 
\begin{itemize}
    \item AC3:
    \item 
\end{itemize}

Para el AC3 de la practica voy a ir reduciendo los dominios de las variables y voy a ir expandiendo  los dominios mas pequeños para buscar la solucion

\section{Ejecrcicios}
\subsection{Ejericio 1} Dada la red inicial de la figura, donde en cada nodo se representa el dominio de la
variable, obténgase la red arco consistente equivalente.
<Falta imagen>

Solucion: <Buscar una bien escrita>

Variables: X1,X2,X3
Restricciones: X1>X2>X3
Paso 1:
Dom(X1)={6,10}
Dom(X2)={5,9}
Dom(X3)={8,15}
Paso 2:
Dom(X1)={6,10}
Dom(X2)={9,9}
Dom(X3)={8,15}
Paso 3:
Dom(X1)={10,10}
Dom(X2)={9,9}
Dom(X3)={8,15}
Paso 4:
Dom(X1)={10,10}
Dom(X2)={9,9}
Dom(X3)={8,8}

\subsection{Ejercicio 2} El sudoku es un claro ejemplo de representaciópn del puzzle como un problema de
satisfacci´on de restricciones. Este problema se publicó en Nueva York en el año 1979
bajo el nombre de “Number Place” y se hizo popular en Japón con el nombre de sudoku
(“Sudji wa dokushin ni kagiru”: “los números deben ser sencillos” o “los números deben
aparecer una vez”).
\medskip
En el problema del sudoku (ver un ejemplo en la figura), hay que rellenar las casillas de
un tablero de 9 x 9 con números del 1 al 9, de forma que no se repita ningún número
en la misma fila, columna, o subcuadro de 3 x 3 que componen el sudoku. Modelar
este problema como un CSP.
<Falta imagen>

\subsection{Ejercicio 3} Cinco casas consecutivas tienen colores diferentes y son habitadas por hombres de
diferentes nacionalidades. Cada uno tiene un animal diferente, una bebida preferida y
fuma una marca determinada. Además, se sabe que:
\begin{itemize}
    \item El noruego vive en la primera casa
    \item La casa de al lado del noruego es azul
    \item El habitante de la tercera casa bebe leche
    \item El inglés vive en la casa roja
    \item El habitante de la casa verde bebe café
    \item El habitante de la casa amarilla fuma Kools
    \item La casa blanca se encuentra justo después de la verde 8. El español tiene un perro
    \item  El ucraniano bebe té
    \item El japonés fuma Cravens
    \item El fumador de Old Golds tiene un caracol
    \item El fumador de Gitanes bebe vino
    \item El vecino del fumador de Chesterfields tiene un reno
    \item El vecino del fumador de Kools tiene un caballo
\end{itemize}

\subsection{Ejercicio 4} Obtener la ordenaci´on de variables seg´un los diferentes m´etodos de ordenaci´on est´aticos
para cada uno de los CSP binarios siguientes:
<Falta imagen>

\subsection{Ejercicio 5} Seis amigos y sus respectivas esposas, se hospedan en el mismo hotel; y todos ellos
salen todos los d´ıas, asistiendo a reuniones de distinto volumen y composici´on. Para
asegurar la variedad en estas salidas, han acordado establecer las siguientes reglas: “Si
Antonio est´a con su mujer, es decir en la misma reuni´on que su mujer, y David con
la suya, y Luis con la se˜nora de Pedro, Enrique debe estar con la se˜nora de Ram´on.
Si Antonio est´a con su mujer y Pedro con la suya, y David con la se˜nora de Enrique,
Ram´on no debe estar con la se˜nora de Luis. Si Enrique y Ram´on y sus mujeres est´an
todos en la misma reuni´on, y Antonio no est´a con la se˜nora de David, Luis no debe
estar con la se˜nora de Pedro. Si Antonio est´a con su mujer y Ram´on con la suya, y
David no est´a con la se˜nora de Enrique, Luis debe estar con la se˜nora de Pedro. Si Luis
est´a con su mujer y Pedro con la suya y Enrique con la se˜nora de Ram´on, Antonio no
debe estar con la se˜nora de David. Si David y Enrique y sus mujeres est´an todos en la
misma reuni´on, y Luis no est´a con la se˜nora de Pedro, Ram´on debe estar con la se˜nora
de Luis”.
\medskip
Modelar dicho problema como un CSP y obtener las posibles combinaciones de reuniones para cada amigo. ¿Es posible que todos los d´ıas haya al menos un matrimonio
cuyos miembros no est´en juntos en la misma reuni´on? Este problema y el siguiente son
originales de Lewis Carroll, en su obra L´ogica Simb´olica.


\subsection{Ejercicio 6} Cinco amigos, Bernardo, Casimir, Luis, Carlos y Marcial, se encuentran cada d´ıa en el
restaurante. Tienen estas reglas, que observan cada vez que comen: “Bernardo toma
sal si y solamente si Casimir toma solo sal o solo mostaza. Bernardo toma mostaza si y
solamente si, o bien Luis no toma sal ni mostaza, o bien Marcial toma ambas. Casimir
toma sal si y solamente si, o Bernardo toma solamente uno de los dos condimentos, o
bien Marcial no toma ninguno de ellos. Toma mostaza si y solamente si Luis o Carlos
toman dos condimentos. Luis toma sal si y solamente si o bien Bernardo no toma
ning´un condimento, o bien Casimir toma ambos. Luis toma mostaza si y solamente
si Carlos o Marcial no toman ni sal ni mostaza. Carlos toma sal si y solamente si
Bernardo o Luis no toman ni sal ni mostaza. Toma mostaza si y solamente si Casimir
o Marcial no toman ni sal, ni mostaza. Marcial toma sal si y solamente si Bernardo o
Carlos toman los dos condimentos. Marcial toma mostaza si y solamente si Casimir o
Luis toman solo un condimento”.
\medskip
Modelar el CSP correspondiente. El problema consiste en descubrir si estas reglas son
compatibles y, en caso de que lo sean, cuales son las combinaciones posibles.

\subsection{Ejercicio 7} Representar mediante un CSP la siguiente informaci´on: Juana, Pepa y Paloma nacieron
y viven en ciudades diferentes (M´alaga, Madrid y Valencia). Adem´as, ninguna vive en
la ciudad donde naci´o. Juana es m´as alta que la que vive en Madrid. Paloma es cu˜nada
de la que vive en Valencia. La que vive en Madrid y la que naci´o en M´alaga tienen
nombres que comienzan por distinta letra. La que naci´o en M´alaga y la que vive ahora
en Valencia tienen nombres que comienzan por la misma letra.
\medskip
¿ D´onde naci´o y vive cada una?

\subsection{Ejercicio 8} Henry Dudeney (1847-1930) era un ingenioso inventor de problemas matem´aticos. Entre sus aportaciones se encuentra un puzzle con cierta complejidad de resoluci´on. Se
trata de un hept´agono donde en cada arista hay que colocar tres n´umeros: uno en cada
v´ertice y otro en el centro de la arista. Hay que colocar los n´umeros del 1 al 14 alrededor
de las aristas del hept´agono de manera que el n´umero de la arista y los dos v´ertices
extremos sumen lo mismo, para todas las aristas. Modelar el CSP correspondiente.
<Falta imagen>

\subsection{Ejercicio 9} Obtener, seg´un las heur´ısticas de ordenaci´on de variables, la mejor ordenaci´on para el
coloreado de los siguientes mapas:


<Falta imagen>


% ---------Bibliografía del Tema------------
\bigskip
\noindent\rule{\textwidth}{0.4pt} 

\section*{Bibliografía del Tema}
\addcontentsline{toc}{section}{Bibliografía del Tema}

\begin{itemize}
    \item
\end{itemize}
\newpage
\chapter{Sistemas Multiagentes}



% ---------Bibliografía del Tema------------
\bigskip
\noindent\rule{\textwidth}{0.4pt} 

\section*{Bibliografía del Tema}
\addcontentsline{toc}{section}{Bibliografía del Tema}

\begin{itemize}
    \item
\end{itemize}
\newpage
\chapter{Web semántica y ontologías}






% ---------Bibliografía del Tema------------
\bigskip
\noindent\rule{\textwidth}{0.4pt} 

\section*{Bibliografía del Tema}
\addcontentsline{toc}{section}{Bibliografía del Tema}

\begin{itemize}
    \item
\end{itemize}
\newpage




% Bibliografia
% \backmatter
\bibliography{Bibliografia}
\bibliographystyle{apacite}
\nocite{*}
\newpage
% Añade la línea al índice manualmente
\addcontentsline{toc}{chapter}{Bibliografía}
\chapter*{Bibliografía}

\begin{itemize}
    \item \textbf{Ana Mas}; \textit{Agentes software y sistemas multiagente. Conceptos, arquitecturas y aplicaciones}; Pearson, 2004.
        
    \item \textbf{Gerhard Weiss (ed.)}; \textit{Multiagent systems}; The MIT Press, 1999.

    \item \textbf{M. Klusch (Ed.)}; \textit{Intelligent Information Agents. Agent-Based Information Discovery and Management}; Springer Verlag, 1999.

    \item \textbf{Russell, S., Norvig, P.}, \textit{Inteligencia Artificial: Un Enfoque Moderno}, Pearson (2004) (4ª Ed).

    \item \textbf{Russell, S. J., \& Norvig, P.} (2020). \textit{Artificial Intelligence: A Modern Approach} (Global Edition).
\end{itemize}

\end{document}